% Part: incompleteness
% Chapter: representability-in-q
% Section: beta-function

\documentclass[../../../include/open-logic-section]{subfiles}

\begin{document}

\olfileid{inc}{req}{bet}
\olsection{బీటా ప్రమేయ ఉపపత్తి}


కూడిక, గుణకారం, $\Char{=}$ అందుబాటులో ఉన్నప్పుడు ఆదిమ పునరావృత్తిని
చేయగలమని చూపడానికి, క్రమాలను నిర్వహించే ప్రమేయాలను అభివృద్ధి చేయాలి.
(ఘాతాంక ప్రమేయం కూడా ఉంటే మన పని సులభమయ్యేది.) ఆదిమ పునరావృత్తి ఉన్నప్పుడు
``$n$-వ ప్రధాన సంఖ్య'' వంటి వాటిని నిర్వచించి, సూటియైన సంకేతీకరణను
ఎంచుకోగలిగాం. కానీ ఇక్కడ ఆదిమ పునరావృత్తి లేదు---నిజానికి కనిష్ఠీకరణను వాడి
ఆదిమ పునరావృత్తి చేయవచ్చని చూపాలనుకుంటున్నాం---కాబట్టి కొంచెం చాకచక్యమైన
పద్ధతి అవసరం.

\begin{lem}
\ollabel{lem:beta}
ప్రతి క్రమం $a_0$,~\dots,~$a_n$కు, ప్రతి $i\le n$కూ
$\beta(d,i)=a_i$ అయ్యే ఒక సంఖ్య~$d$ ఉండేలా ఒక ప్రమేయం~$\beta(d,i)$ ఉంది.
అంతేకాక, ప్రాథమిక ప్రమేయాల నుంచి సంయుక్తం, సక్రమ కనిష్ఠీకరణ మాత్రమే వాడి
$\beta$ను నిర్వచించవచ్చు.
\end{lem}

$d$ అనేది $\tuple{a_0,\dots,a_n}$ క్రమాన్ని సంకేతీకరిస్తోందనీ,
$\beta(d,i)$ దాని $i$-వ మూలకాన్ని ఇస్తోందనీ భావించండి. (ఈ ``సంకేతీకరణ'' మనకు
ఇప్పటికే తెలిసిన ప్రధాన-సంఖ్యా ఘాత సంకేతీకరణను \emph{వాడదని} గమనించండి!)
ఈ ఉపపత్తి చాలా మితమైనదే. క్రమాలను అనుసంధానించవచ్చనీ, మూలకాలను చివర
చేర్చవచ్చనీ అది చెప్పదు; సంయుక్తం, సక్రమ కనిష్ఠీకరణలతో నిర్వచించదగిన
ప్రమేయాలను వాడి $a_0$,~\dots,~$a_n$ల నుంచి~$d$ను \emph{గణించవచ్చని} కూడా
చెప్పదు. ప్రతి క్రమమూ ``సంకేతీకరించబడేలా'' ఒక ``విసంకేతీకరణ'' ప్రమేయం
ఉందని మాత్రమే చెబుతుంది.

$\beta$ సంకేతనం గ్యోడెల్‌ది. మళ్లీ చెప్పాలంటే, పరిమితంగా కనిపించే వనరులతో,
అంటే సంయుక్తం, కనిష్ఠీకరణ మాత్రమే వాడి తగిన~$\beta$ను నిర్వచించడమే ఉపపత్తిని
నిరూపించడంలోని కఠిన భాగం---అయితే కూడిక, గుణకారం, $\Char{=}$లను వాడటానికి
అనుమతి ఉంది. ఈ ఉపపత్తిని నిరూపించడానికి అనేక మార్గాలున్నాయి; వాటిలో ఇప్పటికీ
అత్యంత స్వచ్ఛమైన మార్గాల్లో ఒకటి గ్యోడెల్ అసలు పద్ధతి. అది సున్‌జి సిద్ధాంతం
అనే సంఖ్యా-సిద్ధాంత వాస్తవాన్ని (సాంప్రదాయికంగా ``చైనీయ శేష సిద్ధాంతం'')
వాడింది.

\begin{defn}
రెండు సహజ సంఖ్యలు $a$, $b$ల గరిష్ఠ సామాన్య భాజకం~$1$ అయితే, అప్పుడు మాత్రమే
వాటిని \emph{పరస్పర ప్రధాన సంఖ్యలు} అంటాం; మరో మాటలో, వాటికి మరే సామాన్య
భాజకాలూ ఉండవు.
\end{defn}

\begin{defn}
సహజ సంఖ్యలు $a$, $b$లు \emph{$c$ మాడ్యులో సమానశేషాలు},
$a\equiv b\mod c$, అవుతాయి అంటే $c\mid(a-b)$; అంటే $c$తో భాగించినప్పుడు
$a$, $b$లకు ఒకే శేషం వస్తుంది.
\end{defn}

ఇదే సున్‌జి సిద్ధాంతం:
\begin{thm}
$x_0$,~\dots,~$x_n$లు జతజతగా పరస్పర ప్రధాన సంఖ్యలని అనుకుందాం.
$y_0$,~\dots,~$y_n$లు ఏ సంఖ్యలైనా కావచ్చు. అప్పుడు కింది సమానశేషతలన్నిటినీ
నెరవేర్చే ఒక సంఖ్య~$z$ ఉంటుంది:
\begin{align*}
z & \equiv y_0 \mod x_0 \\
z & \equiv y_1 \mod x_1 \\
& \vdots  \\
z & \equiv y_n \mod x_n.
\end{align*}
\end{thm}

సున్‌జి సిద్ధాంతాన్ని ఇలా వాడతాం: $x_0$,~\dots,~$x_n$లు వరుసగా
$y_0$,~\dots,~$y_n$లకంటే పెద్దవైతే, $\tuple{y_0,\dots,y_n}$ క్రమానికి
$z$ను సంకేతంగా తీసుకోవచ్చు. $y_i$ను తిరిగి పొందడానికి $z$ను~$x_i$తో
భాగించి శేషాన్ని తీసుకుంటే చాలు. ఈ సంకేతీకరణను వాడటానికి
$x_0$,~\dots,~$x_n$లకు తగిన విలువలను కనుగొనాలి.

దీనికి కింది పరిశీలనలు సహాయపడతాయి. $y_0$,~\dots,~$y_n$ ఇచ్చినప్పుడు
\begin{align*}
j &= \max(n, y_0 + 1, \dots, y_n + 1), \\
m &= \lcm(1,\dots,j),
\end{align*}
అని తీసుకుని, తరువాత
\begin{align*}
x_0 & = 1 + m \\
x_1 & = 1 + 2 \cdot m \\
x_2 & = 1 + 3 \cdot m \\
& \vdots  \\
x_n & = 1 + (n+1) \cdot m
\end{align*}
అని తీసుకోండి. అప్పుడు రెండు విషయాలు నిజం:
\begin{enumerate}
\item\ollabel{rel-prime} $x_0,\dots,x_n$లు జతజతగా పరస్పర ప్రధాన సంఖ్యలు.
\item\ollabel{less} ప్రతి $i$కూ $y_i<x_i$.
\end{enumerate}
\olref{rel-prime} నిజమని చూడటానికి, $p$ ఒక ప్రధాన సంఖ్య అనీ,
$p\mid x_i$, $p\mid x_k$ అనీ అనుకుందాం. అప్పుడు
$p\mid1+(i+1)m$, $p\mid1+(k+1)m$. కాబట్టి $p$ వాటి భేదాన్ని భాగిస్తుంది:
\[
(1 + (i+1)m) - (1+ (k+1)m) = (i-k) m.
\]
$p$ అనేది $1+(i+1)m$ను భాగిస్తుంది కనుక అది~$m$ను కూడా భాగించలేదు
(అలా అయితే మొదటి భాగహారంలో శేషం~$1$ వచ్చేది). $p$ అనేది $(i-k)m$ను
భాగిస్తుంది కనుక, $p$ అనేది $i-k$ను భాగిస్తుంది. కానీ $\left|i-k\right|$ గరిష్ఠంగా
$n$; మనం $j\geq n$గా ఎంచుకున్నాం. అందువల్ల $p\mid m$ వస్తుంది; ఇది మళ్లీ
వైరుధ్యం. కాబట్టి $x_i$, $x_k$ రెండింటినీ భాగించే ప్రధాన సంఖ్య ఏదీ లేదు.
\olref{less} చాలా సులభం: $y_i<j\leq m<x_i$.

ఇప్పుడు $\beta$ ప్రమేయ ఉపపత్తిని నిరూపిద్దాం. $0$, ఉత్తరగామి, కూడిక,
గుణకారం, $\Char{=}$, ప్రక్షేపాలు, వాటి నుంచి సంయుక్తం, సక్రమ ప్రమేయాలకు
వర్తింపజేసిన కనిష్ఠీకరణలతో నిర్వచించిన ఏ ప్రమేయాన్నైనా వాడవచ్చని గుర్తుంచుకోండి.
ఒక సంబంధపు లక్షణ ప్రమేయం ఈ విధంగా నిర్వచించదగినదైతే, ఆ సంబంధాన్నీ వాడవచ్చు.
మునుపటిలాగే ఈ సంబంధాలు బూలియన్ సంయోజనలు, పరిమిత పరిమాణీకరణ కింద సంవృతమని
చూపవచ్చు. ఉదాహరణకు:
\begin{align*}
\fn{not}(x) & \defis \Char{=}(x,0)\\
\bmin{x \leq z}{R(x,y)} & \defis \umin{x}{(R(x,y) \lor x = z)}\\
\bexists{x \leq z}{R(x,y)} & \defiff R(\bmin{x \leq z}{R(x,y)}, y)
\end{align*}
ఆదిమ పునరావృత్తి లేకుండానే కింది వాటన్నిటినీ నిర్వచించవచ్చని కూడా చూపవచ్చు:
\begin{enumerate}
\item జతీకరణ ప్రమేయం,
$J(x,y)=\frac{1}{2}[(x+y)(x+y+1)]+x$;
% maybe explain more what is going on here, a bit confusing.
\item ప్రక్షేప ప్రమేయాలు
\begin{align*}
K(z) & = \bmin{x \leq z}{\bexists{y \leq z}{z = J(x,y)}},\\
L(z) & = \bmin{y \leq z}{\bexists{x \leq z}{z = J(x,y)}};
\end{align*}
\item చిన్నది-కంటే సంబంధం $x<y$;
\item భాగించు సంబంధం $x\mid y$;
% \item $x \tsub y$
% \item $\fn{Prime}(x)$
% \item Assuming $p$ is prime, the relation ``$x$ is a power of $p$'':
% \[
% \bforall{y \leq x}{(y \mid x \lif y = 1 \lor y = x)}.
% \]
\item $y$ను~$x$తో భాగించినప్పుడు వచ్చే శేషాన్ని ఇచ్చే
  ప్రమేయం~$\fn{rem}(x,y)$.
\end{enumerate}
ఇప్పుడు ఇలా నిర్వచించండి:
\begin{align*}
\beta^*(d_0,d_1,i) & = \fn{rem}(1+(i+1) d_1,d_0) \text{ మరియు}\\
\beta(d,i) & = \beta^*(K(d),L(d),i).
\end{align*}
మనకు కావలసిన ప్రమేయం ఇదే. పైన చెప్పిన $a_0,\dots,a_n$లు ఇచ్చినప్పుడు
\[
j = \max(n,a_0+1,\dots,a_n+1),
\]
అని, $d_1=\lcm(1,\dots,j)$ అని తీసుకోండి. పై \olref{rel-prime} వల్ల
$1+d_1$, $1+2d_1$,~\dots,~$1+(n+1)d_1$లు జతజతగా పరస్పర ప్రధాన సంఖ్యలని,
\olref{less} వల్ల అవి వరుసగా $a_0,\dots,a_n$లకంటే పెద్దవని తెలుసు. సున్‌జి
సిద్ధాంతం వల్ల ప్రతి~$i$కూ
\[
d_0 \equiv a_i \mod (1+(i+1)d_1)
\]
అయ్యే ఒక విలువ~$d_0$ ఉంటుంది. అందువల్ల ($d_1$ అనేది~$a_i$కంటే పెద్దది కనుక)
\[
a_i = \fn{rem}(1+(i+1)d_1,d_0).
\]
$d=J(d_0,d_1)$గా తీసుకోండి. అప్పుడు ప్రతి $i\le n$కూ
\begin{align*}
\beta(d,i) & =  \beta^*(d_0,d_1,i) \\
& =  \fn{rem}(1+(i+1) d_1,d_0) \\
& =  a_i
\end{align*}
వస్తుంది; మనకు కావలసింది ఇదే. దీంతో $\beta$-ప్రమేయ ఉపపత్తి నిరూపణ
పూర్తయింది.

\begin{prob}
  ఆదిమ పునరావృత్తి లేకుండా $x<y$, $x\mid y$ సంబంధాలు,
  $\fn{rem}(x,y)$ ప్రమేయాన్ని నిర్వచించవచ్చని చూపండి. $0$, ఉత్తరగామి,
  కూడిక, గుణకారం, $\Char{=}$, ప్రక్షేపాలు, పరిమిత కనిష్ఠీకరణ, పరిమిత
  పరిమాణీకరణలను వాడవచ్చు.
\end{prob}

\end{document}
